\section{Related Work}
\label{sec:related}

\textbf{Interactive environments for coding.} 
Most coding benchmarks (e.g.\,SQL - Spider \cite{yu-etal-2018-spider}, KaggleDBQA \cite{lee-etal-2021-kaggledbqa}; Bash - NLC2CMD \cite{pmlr-v133-agarwal21b}, NL2Bash \cite{lin-etal-2018-nl2bash}; Python - HumanEval \cite{chen2021evaluating}, APPS \cite{hendrycksapps2021}, MBPP \cite{austin2021program}, CodeXGLUE \cite{lu2021codexglue}, CodeNet \cite{puri2021codenet}) frame the coding problem as a sequence transduction problem (from instruction to code), rather than an interactive decision making problem with an execution environment. 
Attempts have been made to simulate interaction by developing conversational, dialogue-style \cite{yu-etal-2019-sparc, yu2019cosql}, multi-step problem solving \cite{nijkamp2022codegen} datasets, which involve pre-annotated human-designed queries. The work closest to \benchmark{} has been recent explorations of Python Jupyter Notebooks as a natural choice for interactive coding \cite{huang-etal-2022-execution, Lai2022DS1000, yin2022natural}. However, task data and settings often constrain allowed actions to a closed domain of code and libraries~\cite{Lai2022DS1000, yin2022natural}, use evaluation procedures or metrics that may not generalize~\cite{huang-etal-2022-execution}, require human-in-the-loop participation (i.e. create task contexts, write problems, evaluate execution per task instance)~\cite{Lai2022DS1000}, or are Python-exclusive~\cite{huang-etal-2022-execution, Lai2022DS1000, yin2022natural, wang2023executionbased}.
% In addition, such benchmarks cannot be adapted to other programming languages or task settings in a straightforward manner.
\benchmark{} provides a more general purpose foundation defining interactive coding tasks that enables easy construction of diverse task settings, can have any programming language(s) as the action space, and has automatic, execution-based evaluation.

\textbf{Execution-based evaluation for coding.} Evaluation for NL-to-code generation models has recently shifted away from surface form similarity metrics (BLEU~\cite{Papineni2002BleuAM, Agashe2019JuICeAL}, ROUGE~\cite{lin-2004-rouge}, Exact Match) towards execution oriented ratings (unit tests~\cite{austin2021program, chen2021evaluating, huang-etal-2022-execution, Lai2022DS1000, hendrycksapps2021}, output matching~\cite{dong2016language, huang-etal-2022-execution, zhong2017seq2sql}).
The rigidity of surface form analysis overlooks code syntax features, ignores execution effect, or over-penalizes alternative solutions~\cite{zhou2023codebertscore},
On the contrary, execution-based assessment is a more thorough and comprehensive score of code functionality~\cite{hendrycksapps2021} and is a more natural fit for open-domain program usage that does not constrain code generation to a subset of the language space~\cite{wang2023executionbased}.
However, for newer benchmarks and datasets that put forth task definitions incorporating execution-based evaluation (APPS~\cite{hendrycksapps2021}, ExeDS~\cite{huang-etal-2022-execution},  ODEX~\cite{wang2023executionbased}), the fundamental code generation task (Context + Code $\rightarrow$ Execution $\rightarrow$ Score) is still devoid of interaction.
% In addition, the proposal of new metrics borne from different task formulations may make it difficult to compare models across all settings.
\benchmark{} combines execution-based evaluation with flexible task construction, enabling more diverse problem-solving paradigms within a unified coding task formulation.
\benchmark{}'s use of virtual containers as execution sandboxes protect against harmful actions and allow for advanced evaluation criteria beyond the aforementioned ones.
% surface form, unit tests, or execution output matching.
% An example is the \benchmark{}-Bash environment, which includes file system modification in its task completion function, a capability not easily achievable with previous evaluation methods.

\textbf{Methods for interactive or execution-based coding.} The value of generative code models and interactive problem solving has motivated a recent proliferation of work to augment reasoning capabilities' of existing language models~\cite{yao2023react, shinn2023reflexion, wang2023planandsolve, yao2023tree, zhang2023planning, chen2023teaching} or propose new modeling techniques to tackle coding as a sequential decision making and reasoning tasks~\cite{bunel2018leveraging, Chen2018ExecutionGuidedNP, ellis2019write, alphacode, chen2022codet, le2022coderl}, of which evaluation is unit test based.
Approaches that leverage execution typically use re-ranking~\cite{zhang2022coder, ni2023lever, yin-neubig-2019-reranking, zeng2022nbest} or majority vote~\cite{Chen2018ExecutionGuidedNP, alphacode, shi2022natural} to decide on a final prediction.
Additional work also explores incorporating human-in-the-loop~\cite{chen2023improving, lahiri2022interactive}, compilers~\cite{wang2022compilable}, and text~\cite{wang2023leti,zhang2023selfedit} feedback.
A common thread among these contributions is that 1) the task setting can only provide the investigated form of feedback and 2) sought-after capabilities are exemplified by strong performance on favorably curated tasks and datasets, rendering comparisons across benchmarks tedious.
\benchmark{} has the potential to standardize the evaluation of these methods because 1) the interactive coding task is a conglomeration of many interesting interaction, reasoning, and decision-making challenges and 2) \benchmark{}'s task construction makes it possible to incorporate a wide variety of sources of feedback.

% \textbf{Other language-based interaction} Text games, web navigation, etc. but natural langauge, not programming language as action. \sy{less important, hold for now} 

% \textbf{Interactive coding benchmarks}
% % existing benchmarks
% %As models become increasingly efficient, developing more challenging and general code generation benchmarks has become the norm. 
% Most coding benchmarks (e.g.\,SQL \cite{yu-etal-2018-spider,lee-etal-2021-kaggledbqa}, bash \cite{pmlr-v133-agarwal21b, lin-etal-2018-nl2bash}, Python \sy{explicitly write down names}) set up a coding problem as an instruction and a .

% In the absence of an interactive coding benchmark that can cater to more than just Python, we develop \benchmark{}, which currently provides interactive environments for SQL, Bash, and Python and demonstrate that it can be extended to other exciting code-based tasks like CTFs.



% \textbf{Interactive or execution-based coding methods} Several past and recent works employ variations of test-driven development
% or code debugging practices. AlphaCode~\citep{li2022competition}
% evaluates a set of generations on hidden test cases. CodeT~\citep{chen2022codet}
% uses self-generated unit tests that are used to score generated function
% implementations.
% Self-Debugging~\citep{chen2023teaching} employs a debugging component that is used
% to improve existing implementations given feedback from a code
% execution environment. CodeRL~\citep{le2022coderl} sets the problem
% in an RL framework using an actor-critic setup to debug programs given
% feedback from an execution environment. AlphaCode, Self-Debugging and CodeRL
% are effective in fixing less-complex program bugs, but they rely
% upon ground truth test cases that invalidate pass@1 eligibility,
% and do not use self-reflection to bridge the gap between error
% identification and implementation improvement. CodeT does not access
% hidden test cases but does not implement a self-learning step to
% improve code writing.




% The standard approach to generating code involves models generating a blob of code given a natural language intent and are typically tested via hand-designed test cases \cite{chen2021evaluating, alphacode,hendrycksapps2021}.
% % multi-turn: decompose
% Progressive generation of subprograms through multi-turn expert-generated prompts \cite{nijkamp2022codegen} demonstrated that Language Models could benefit from the systematic multi-step approach to code generation but offers no scope for error-recovery in case of incorrectly generated previous turns. 
% Recently, language models have emerged as agents capable of observing, interacting, and receiving iterative feedback from external entities and the environment \cite{wang2023interactive}. 

% - Planning \cite{zhang2023planning}

% - Human-in-the-loop \cite{chen2023improving, lahiri2022interactive}

% - Compiler feedback \cite{wang2022compilable}

% - Textual feedback \cite{wang2023leti}



% \textbf{Language-based Interaction} \sy{less important} 
% % \textbf{Text Games}
% % CALM: Similar challenge of determining both admissible and promising actions when the action space is extremely large (essentially infinite)
